\documentclass[UTF8, a4paper]{ctexart}

% ===== 【新增】解决打印字体发浅的代码 =====
\setCJKmainfont{SimSun}[AutoFakeBold=1.2] % 强制使用宋体，并增加1.2倍的伪粗体字重，适配物理打印
\PassOptionsToPackage{text}{xeCJK}        % 优化汉字边缘的渲染平滑度，防止打印发虚
% ==========================================

\usepackage{amsmath}
\usepackage{amssymb}
\usepackage{fancyhdr}
\usepackage{titlesec}

% --- 页码样式设置 ---
\pagestyle{fancy}
\fancyhf{}
\fancyfoot[C]{\thepage}
\renewcommand{\headrulewidth}{0pt}
% --- 页码样式设置结束 ---

% --- 关闭章节自动编号 ---
\setcounter{secnumdepth}{0}

% --- 【核心修改】重写 \part 的排版，让它完美居中 ---
\titleformat{\part}[display]
  {\normalfont\huge\bfseries\centering} % 字体巨大、居中
  {\partname\ \thepart}                 % 显示“第一部分”
  {20pt}                                % 标题和编号之间的距离
  {\Huge\bfseries}                      % 正文部分标题
\titlespacing*{\part}{0pt}{50pt}{50pt}  % 上下各留50pt的留白，强制居中
% --- 【核心修改】结束 ---
\newcommand{\dd}{\mathrm{d}}
\usepackage{amsmath}
\makeatletter
\let\old@eqnnum=\@eqnnum
\renewcommand{\@eqnnum}{\relax}
\makeatother
\begin{document}

% --- 生成目录 ---
{
\pagestyle{empty}
\tableofcontents
\newpage
\setcounter{page}{1}
}
% --- 目录结束 ---

% ==============================
% 第一部分 物质、空间与时间
% ==============================
\part{物质、空间与时间}
\include{1-1}
\include{1-2}

% ==============================
% 第二部分 质点运动学
% ==============================
\part{质点运动学}
\include{2-1}

% ==============================
% 第三部分 牛顿力学
% ==============================
\part{牛顿力学}
\include{3-1}
\include{3-2}

% ==============================
% 第四部分 振动与波
% ==============================
\part{振动与波}
\include{4-1}
\include{4-2}
\include{4-3}

% ==============================
% 第五部分 电磁学
% ==============================
\part{电磁学}
\include{5-1}
\include{5-2}
\include{5-3}
\include{5-4}
\include{5-5}

% ==============================
% 第六部分 几何光学
% ==============================
\part{几何光学}
\include{6-1}

% ==============================
% 第七部分 热力学
% ==============================
\part{热力学}
\include{7-1}
\include{7-2}
\include{7-3}
\include{7-4}

% ==============================
% 第八部分 分析力学
% ==============================
\part{分析力学}
\include{8-1}
\include{8-2}
\include{8-3}



\part{哈密顿力学}
\include{9-1}
\include{9-2}
\include{9-3}

\end{document}